% Anidamiento o alcance: un dominio, modelo o teoría dentro de otro. El contorno codifica el
% estatus: sólido = actual, rayado = virtual. El encastre espacial ES la inclusión.
\documentclass[tikz,border=12pt]{standalone}
\usepackage{diagrama}
\begin{document}
\begin{tikzpicture}[every node/.append style={font=\px{28}}]
  \node[titulo=44] at (3.6,5.0) {Título del anidamiento};
  \node[caja, minimum width=7.2cm, minimum height=4.4cm] at (2.25,2.5) {};
  \node[font=\bfseries\px{30}, anchor=north] at (2.25,3.75) {Teoría Tj};
  \node[caja virtual, minimum width=3.4cm, minimum height=1.9cm] (n1) at (2.25,2.2) {Modelo\\contraparte};
  \node[caja, minimum width=2.4cm, minimum height=1cm] (n3) at (6.6,2.2) {Teoría Tk};
  \draw[virtual] (n1.east) -- node[rotulo=20] {\textit{Isomorfo}} (n3.west);   % e1
\end{tikzpicture}
\end{document}
